% Source: https://artofproblemsolving.com/community/c3h1371579_how_many_unit_squares_does_this_line_pass_through
% Stable identity: aops-topic-1371579-q01
Place the rectangle with corners $(0,0)$ and $(7,5)$; the diagonal is $y=\dfrac{5}{7}x$.

\begin{asy}
usepackage("amsmath");
size(6cm);
defaultpen(fontsize(10pt));
for(int i=0;i<7;++i) for(int j=0;j<5;++j) if(5*i/7<j+1 && 5*(i+1)/7>j) fill(shift(i,j)*unitsquare,lightblue);
for(int i=0;i<=7;++i) draw((i,0)--(i,5),gray(0.5));
for(int j=0;j<=5;++j) draw((0,j)--(7,j),gray(0.5));
draw((0,0)--(7,0)--(7,5)--(0,5)--cycle,linewidth(0.9));
draw((0,0)--(7,5),red+linewidth(0.9));
label("$7$",(3.5,0),S); label("$5$",(7,2.5),E);
label("\mbox{}",(min(currentpicture,user=true)+max(currentpicture,user=true))/2,black); // keeps the text after the figure black when LaTeX typesets the labels inline
\end{asy}

Travel along the diagonal from $(0,0)$. It starts in one square, and it enters a new square each time it crosses an interior grid line: the $6$ vertical lines $x=1,\dots,6$ and the $4$ horizontal lines $y=1,\dots,4$. It never crosses two grid lines at once, because that would happen at a lattice point $(a,b)$ with $0<a<7$ and $b=\dfrac{5a}{7}$ an integer, which is impossible as $\gcd(5,7)=1$. So the number of squares is
\[
1+6+4=\boxed{11}.
\]
(In general, an $m\times n$ diagonal passes through $m+n-\gcd(m,n)$ squares.)

Column by column (each column is $1$ wide, and the diagonal rises $\dfrac{5}{7}$ in it), the diagonal meets $2$ squares in a column that contains a crossing of a horizontal line and $1$ square otherwise. It crosses $4$ horizontal lines in $7$ columns, so $4$ columns contain $2$ squares and $3$ columns contain $1$: $4\cdot2+3=11$. Turning the rectangle the other way (slope $\dfrac{7}{5}>1$, as in the question's attempt), the diagonal crosses $6$ lines in $5$ columns, so one column must contain two crossings and $3$ squares: in the middle column the diagonal rises from $\dfrac{14}{5}$ to $\dfrac{21}{5}$, crossing $y=3$ and $y=4$. That column has $3$ squares and the others $2$: $3+4\cdot2=11$ again.
